\documentclass[12pt]{article}

% ==== core math ====
\usepackage{amsmath,amssymb,amsthm,amsfonts}
% ==== fonts (bold typewriter for \paragraph headings with \texttt) ====
\usepackage{bold-extra}
% ==== diagrams ====
\usepackage{tikz}
\usepackage{tikz-cd}
% ==== graphics ====
\usepackage{graphicx}
% ==== tables ====
\usepackage{booktabs}
% keep floats after their point of definition (Table 1 after its Definition)
\usepackage{flafter}
% ==== page layout ====
\usepackage[margin=1in]{geometry}
% ==== references (hyperref before cleveref; aliascnt before cleveref) ====
\usepackage{hyperref}
\usepackage{aliascnt}
\usepackage{cleveref}
% ==== GrokRxiv sidebar ====
\usepackage{xcolor}

\hypersetup{
  colorlinks=true,
  linkcolor=blue!45!black,
  citecolor=green!35!black,
  urlcolor=blue!55!black
}

% Absorb the occasional stubborn line rather than overflow the margin.
\emergencystretch=2em

% ==== topological-phases-of-matter : canonical macros (v1, 2026-07-14) ====
% Requires amsmath, amssymb, amsfonts (mathfrak, mathbf, mathcal all standard).

% --- condensation / probes ---
\newcommand{\cond}[1]{\underline{#1}}          % condensation X |-> \cond{X}
\newcommand{\cg}[1]{\underline{#1}}            % condensed group (alias, for readability)
\newcommand{\Cont}{\operatorname{Cont}}        % Cont(S,X)
\newcommand{\Shape}{\operatorname{Shape}}      % shape / homotopy type

% --- interaction space & locality ---
\newcommand{\Int}{\mathcal{I}}                 % \Int_{d,G} interaction space
\newcommand{\Ffun}{F}                          % F-function (interaction decay)
\newcommand{\BF}{\mathcal{B}_{F}}              % interaction Banach space for F

% --- moduli stacks ---
\newcommand{\Ham}{\mathfrak{Ham}}              % \Ham_{d,G}
\newcommand{\Gap}{\mathfrak{Gap}}              % \Gap_{d,G}
\newcommand{\Phase}{\mathfrak{Phase}}          % \Phase_{d,G}
\newcommand{\Phases}{\operatorname{Phases}}    % \Phases_{d,G} = pi_0 Shape Phase
\newcommand{\Weq}{\mathcal{W}}                 % class of equivalences to invert
\newcommand{\st}{\mathrm{st}}                  % stabilization superscript
\newcommand{\hfix}[1]{h#1}                     % homotopy-fixed-point superscript, e.g. \Phase_d^{\hfix{\cg{G}}}
\newcommand{\quotstack}[2]{[#1/#2]}            % action stack [Y/\cg{G}]

% --- gap ---
\newcommand{\gapof}[1]{\operatorname{gap}(#1)} % gap(H)
\newcommand{\gapbound}{\Delta}                  % uniform gap bound

% --- stacking / spectrum ---
% \boxtimes is built in (stacking product)
\newcommand{\IPcond}{\mathbf{IP}^{\mathrm{cond}}}  % invertible condensed phase spectrum

% --- transitions / discriminant / charge ---
\newcommand{\Sig}{\Sigma}                       % \Sig_f gapless discriminant
\newcommand{\rc}{\partial\nu}                   % relative transition charge in E^{q+1}(B,U_f)

% --- observables & solid K-theory ---
\newcommand{\Kop}{K_{\mathrm{op}}}             % operator / topological K-theory
\newcommand{\Kalg}{K_{\mathrm{alg}}}           % algebraic K-theory
\newcommand{\Solid}{\operatorname{Solid}}      % solidification functor

% --- disorder ---
\newcommand{\dis}{\Omega}                       % disorder hull \Omega = Q^{Z^d} (Q = local-configuration alphabet; NOT the F-function)
% ==== end canonical macros ====

% --- paper-local shorthands (do not conflict with canonical block) ---
\newcommand{\Pic}{\operatorname{Pic}}
\newcommand{\Groth}{K}                       % Grothendieck group completion of a monoid
\newcommand{\Sph}{\mathbb{S}}
\newcommand{\IZ}{I_{\mathbb{Z}}}             % Anderson dual of the sphere
\newcommand{\IFH}{\mathbf{I}_{\mathrm{FH}}}  % Freed--Hopkins invertible-TQFT spectrum
\newcommand{\Zp}{\mathbb{Z}}
\newcommand{\Rp}{\mathbb{R}}
\newcommand{\Tt}{\mathbb{T}}
\newcommand{\MTH}{\mathit{MTH}}
\newcommand{\triv}{\mathbf{1}}               % trivial (product-state) phase
\newcommand{\inv}[1]{\Phase_{#1}^{\times}}   % invertible sector
\DeclareMathOperator{\coker}{coker}
\DeclareMathOperator{\im}{im}
\DeclareMathOperator{\Map}{Map}
\DeclareMathOperator{\colim}{colim}
\DeclareMathOperator{\KO}{KO}

% ==== theorem environments ====
% aliascnt lets Propositions/Lemmas/etc. share the Theorem counter while keeping
% their own names, so cleveref refers to each by its true kind (not "Theorem").
\newtheorem{theorem}{Theorem}[section]

\newaliascnt{proposition}{theorem}
\newtheorem{proposition}[proposition]{Proposition}
\aliascntresetthe{proposition}

\newaliascnt{lemma}{theorem}
\newtheorem{lemma}[lemma]{Lemma}
\aliascntresetthe{lemma}

\newaliascnt{corollary}{theorem}
\newtheorem{corollary}[corollary]{Corollary}
\aliascntresetthe{corollary}

\theoremstyle{definition}
\newaliascnt{definition}{theorem}
\newtheorem{definition}[definition]{Definition}
\aliascntresetthe{definition}

\newaliascnt{example}{theorem}
\newtheorem{example}[example]{Example}
\aliascntresetthe{example}

\theoremstyle{remark}
\newaliascnt{remark}{theorem}
\newtheorem{remark}[remark]{Remark}
\aliascntresetthe{remark}

% Numbered Conjectures carry STABLE series IDs (IV-1..IV-5); unnumbered env, ID in title.
\theoremstyle{plain}
\newtheorem*{conjecture}{Conjecture}

% cleveref names for the aliascnt-backed environments
\crefname{proposition}{Proposition}{Propositions}
\Crefname{proposition}{Proposition}{Propositions}
\crefname{lemma}{Lemma}{Lemmas}
\Crefname{lemma}{Lemma}{Lemmas}
\crefname{corollary}{Corollary}{Corollaries}
\Crefname{corollary}{Corollary}{Corollaries}
\crefname{definition}{Definition}{Definitions}
\Crefname{definition}{Definition}{Definitions}
\crefname{example}{Example}{Examples}
\Crefname{example}{Example}{Examples}

% ==== GrokRxiv DOI sidebar ====
\definecolor{grokgray}{RGB}{110,110,110}
\AddToHook{shipout/background}{%
  \ifnum\value{page}=1
    \begin{tikzpicture}[remember picture, overlay]
      \node[
        rotate=90,
        anchor=south,
        font=\Large\sffamily\bfseries\color{grokgray},
        inner sep=0pt
      ] at ([xshift=38pt, yshift=0.52\paperheight]current page.south west)
      {GrokRxiv:2026.07.lattice-eft-equivalence\quad
       [\,math.AT\,]\quad
       14 Jul 2026};
    \end{tikzpicture}
  \fi
}

\title{From Lattice Models to Effective Field Theories:\\
Stabilization and the Invertible Condensed Phase Spectrum}

\author{Matthew Long \\
\textit{The YonedaAI Collaboration} \\
\textit{YonedaAI Research Collective} \\
Chicago, IL \\
\texttt{matthew@yonedaai.com} $\cdot$ \url{https://yonedaai.com}}

\date{14 July 2026}

\begin{document}
\maketitle

\begin{abstract}
We study the fourth analytic input to the condensed-mathematics program for
topological phases of matter: the relationship between microscopic lattice
systems and their effective field theories. Over the stabilized phase
$\infty$-groupoid $\Phase_{d,G}$ built in Part~III we make the stacking
operation $\boxtimes$ explicit as a symmetric-monoidal structure and record what
follows formally. On components, stacking makes the set of phases a commutative
monoid whose unit is the trivial product state and whose invertible elements are
exactly the short-range-entangled (SRE) phases. Group completion of this monoid
is governed by a universal property that we prove, and its explicit presentation
isolates a \emph{stabilization element} whose physical content is the addition of
trivial ancillas. Restricting to the invertible sector and invoking the
recognition principle for grouplike $E_\infty$-spaces yields a connective
spectrum, the invertible condensed phase spectrum $\IPcond_{d,G}$; Kubota's
operator-algebraic $\Omega$-spectrum is a rigorous carrier for its homotopy
groups, and Aoki's solidification bridge supplies the analytic completion that a
condensed refinement requires. Against this we set the effective-field-theory
(bordism / Anderson-dual) classification of Freed--Hopkins. The organizing claim
of the program, that the microscopic and field-theoretic classifications
agree, is stated as a comparison conjecture at three increasingly strong levels:
on $\pi_0$ of phase sets, on spectra, and on families over condensed probes. We
do not prove it. We assemble the evidence that constrains it: in $d=1$ Ogata's
operator-algebraic index is a complete invariant and the $\pi_0$ comparison is a
theorem, and for free fermions the tenfold-way $K$-theory of Kitaev realizes it,
with the $8$-fold real and $2$-fold complex Bott periodicities verified by
machine. Renormalization is framed as a filtered / pro-completion, and we
conjecture that the lattice-to-EFT passage is computed by solidification, that
solidification is compatible with group completion, and that the relative
transition charge is the boundary map of the cofiber sequence for $\IPcond$. Two
honest limitations bound the discussion: the equivalence is expected only under
short-range-entanglement hypotheses, and noninvertible topological order lies
outside the invertible spectrum entirely. Accompanying Haskell code implements
finitely presented commutative monoids, their group completion, and the
tenfold-way table, with QuickCheck tests of the monoid laws, the universal
property, and Bott periodicity.
\end{abstract}

\tableofcontents

\section{Introduction}
\label{sec:intro}

This is the fourth paper of a six-part series that reorganizes the theory of
topological phases of matter inside condensed mathematics. The program builds a
single tower,
\begin{equation}
\label{eq:tower}
\text{local interactions}
\ \leadsto\ \Ham_{d,G}
\ \leadsto\ \Gap_{d,G}
\ \leadsto\ \Phase_{d,G}
\ \leadsto\ \IPcond_{d,G},
\end{equation}
whose floors are, respectively: the condensed moduli stack of $G$-symmetric
quasi-local Hamiltonians (Part~I,~\cite{paperLocality}); its uniformly gapped
substack, with the observable side and its solid $K$-theory installed in
Part~II,~\cite{paperPositivity}; the stabilized phase $\infty$-groupoid whose
existence and robustness rest on the gap-stability results of Part~III,
\cite{paperGap}; and, at the top, a connective spectrum assembled from the
invertible sector. The present paper is about the top two floors of
\eqref{eq:tower} and, in particular, about the bridge from a lattice system to
its effective field theory (EFT).

\subsection{The lattice--EFT problem}

A gapped lattice Hamiltonian is a microscopic object: finite-dimensional Hilbert
spaces on sites, a local interaction, an energy gap above the ground state. Its
low-energy physics is described, when the system is short-range entangled, by an
\emph{invertible} topological quantum field theory: an EFT with a
one-dimensional state space on every closed spatial slice
\cite{freed-sre,freed-hopkins}. Two classification schemes then coexist:
\begin{itemize}
\item the \emph{microscopic} one, which sorts lattice systems into phases by
  adiabatic / finite-depth equivalence and stabilization, producing the set
  $\Phases_{d,G}=\pi_0\,\Shape(\Phase_{d,G})$; and
\item the \emph{field-theoretic} one, which sorts invertible TQFTs into
  deformation classes, producing (under reflection positivity) a bordism /
  Anderson-dual invariant computed by a spectrum $\IFH$~%
  \cite{freed-hopkins,kapustin-cobordism,gaiotto-jf,xiong-minimalist}.
\end{itemize}
The fourth item on the program's list of analytic obligations
(\cite{paperLocality}, \S1.8) is to control the relationship between these two.
We take ``equivalence between microscopic lattice systems and effective field
theories'' to mean the assertion that the two classifications agree, and we take
seriously the fact that this is, at present, \emph{open}. It is the organizing
conjecture of the paper, not a theorem we are going to prove.

\subsection{What is proved, and what is conjectured}

The paper keeps a strict line between the two. Provable today, from cited
present-day results, are the algebraic and homotopical facts that make the top of
the tower well-posed:
\begin{itemize}
\item stacking $\boxtimes$ endows $\Phases_{d,G}$ with the structure of a
  commutative monoid (\Cref{prop:monoid}); the trivial phase is the unit and the
  SRE phases are the invertible elements;
\item the Grothendieck group completion of a commutative monoid exists and is
  characterized by a universal property (Theorem~IV-A, \Cref{thm:groth}), whose
  explicit form exposes the ``stabilization element'' behind physical
  stabilization by~%
  ancillas (\Cref{prop:stable});
\item the invertible sector is grouplike, hence---by the recognition principle
  for grouplike $E_\infty$-spaces used in this setting by~%
  \cite{beaudry-etal,kubota-omega-spectrum}---the infinite-loop space of a
  connective spectrum, for whose homotopy groups Kubota constructs a rigorous
  $\Omega$-spectrum carrier (Theorem~IV-B, \Cref{thm:kubota});
\item the free-fermion tenfold way is $8$-fold ($2$-fold) Bott periodic in the
  real (complex) classes (\Cref{prop:bott}); and
\item the $2$D Hall conductance of an SRE state is an integer-valued phase
  invariant, locally computable, with a higher Berry generalization for families
  (Theorem~IV-C, \Cref{thm:hall}).
\end{itemize}
Everything past this is stated as a numbered \emph{Conjecture} with a stable
identifier of the form IV-$n$. Five appear: the condensed / solid refinement of
the spectrum (IV-1), the lattice--EFT comparison equivalence (IV-2, the
organizing conjecture), the identification of the relative transition charge with
a spectral boundary map (IV-3), the compatibility of solidification with group
completion (IV-4), and the existence of a renormalization functor computing the
EFT limit (IV-5).

\subsection{Three levels of equivalence}
\label{subsec:levels}

The word ``equivalence'' hides a hierarchy that we insist on making visible.
There are (at least) three honestly different statements one might mean:
\begin{enumerate}
\item[(L0)] \textbf{Sets of phases.} A bijection $\Phases_{d,G}\cong
  \pi_0\IFH$ of abelian groups: the microscopic and field-theoretic invariants
  classify the same phases.
\item[(L1)] \textbf{Spectra.} An equivalence $\IPcond_{d,G}\simeq\IFH$ (after a
  specified completion) of connective spectra: not only $\pi_0$ but the entire
  homotopy type---higher pumps and families---match.
\item[(L2)] \textbf{Families.} A natural equivalence over condensed probes $S$,
  so that the comparison respects continuous families, profinite disorder hulls
  $\dis=Q^{\mathbb{Z}^d}$, and symmetry data simultaneously.
\end{enumerate}
These are strictly increasing in strength: (L2)$\Rightarrow$(L1)$\Rightarrow$(L0).
Existing anchors settle (L0) in favourable cases (Ogata in $d=1$; free fermions
in every $d$) and give rigorous \emph{targets} for (L1) (Kubota's
$\Omega$-spectrum). Nothing in the literature establishes (L2), and the
condensed refinement it would require is exactly Conjecture IV-1. Keeping the
levels apart is not pedantry: a proof at (L0) leaves (L1) and (L2) untouched, and
most of what one wants the tower \eqref{eq:tower} for (parametrized invariants,
descent along disorder hulls, the transition calculus) lives at (L1) and (L2).

\subsection{Relation to companion papers}
\label{subsec:companions}

This paper consumes the outputs of Parts~I--III and feeds Parts~V--VI.

\emph{Part~I (Condensed Locality}~\cite{paperLocality}\emph{)} makes the
interaction space a Banach space $\BF$ and the Heisenberg dynamics a morphism of
condensed objects; it is the reason a ``family of systems'' over a probe $S$ is a
well-defined thing to stack. \emph{Part~II (Positivity and $C^*$-Norms}
\cite{paperPositivity}\emph{)} attaches the condensed observable algebra
$\cond{A}$ and, through Aoki's theorem $\Kop(A)\simeq\Solid(\Kalg(\cond{A}))$,
the solid $K$-theory that we use here as the analytic completion step in the EFT
passage (\Cref{sec:rg}). \emph{Part~III (The Uniformly Gapped Substack}
\cite{paperGap}\emph{)} is what makes $\pi_0$ of the stabilized stack a robust
invariant in the first place; without the stability of the gap under quasi-local
perturbations, the monoid $\Phases_{d,G}$ would not be well-defined, and the
undecidability wall of Cubitt--P\'erez-Garc\'ia--Wolf~\cite{cpw-undecidable}
forbids us from pretending otherwise. \emph{Part~V (Physical Realizability}
\cite{paperRealizability}\emph{)} takes over exactly where our comparison
conjecture ends: it asks which classes in $\IPcond_{d,G}$ are realized by honest
lattice families, a question sharpened by the Kapustin--Fidkowski no-go
\cite{kapustin-fidkowski}. \emph{Part~VI (Synthesis}~\cite{paperSynthesis}\emph{)}
assembles all five modules and states the global Master Conjecture; our
Conjecture IV-2 is its top floor. We import notation verbatim from the series'
canonical conventions and do not restate it.

\subsection{Outline}
\Cref{sec:recollections} recalls the stabilized stack and fixes the stacking
product. \Cref{sec:monoidal,sec:groupcompletion} carry the provable algebra:
the symmetric-monoidal structure and its group completion.
\Cref{sec:spectrum} assembles the invertible condensed phase spectrum.
\Cref{sec:rg} formulates renormalization as a pro-completion and states the
solidification conjectures. \Cref{sec:comparison} states the comparison
conjecture at all three levels, and \Cref{sec:evidence} marshals the evidence:
$d=1$ completeness and free-fermion Bott periodicity. \Cref{sec:invariants}
treats concrete lattice invariants, the SSH example, and transitions.
\Cref{sec:haskell} describes the accompanying machine verification, and
\Cref{sec:discussion,sec:conclusion} discuss limitations and conclude.

\section{The stabilized phase stack and its stacking product}
\label{sec:recollections}

We recall only what is needed and refer to~\cite{paperGap} for the analytic
construction. Fix a spatial dimension $d$, a lattice $L$, on-site Hilbert spaces,
an $F$-function locality class, and a symmetry group $G$ with condensation
$\cg{G}$.

\begin{definition}[the tower,~\cite{paperLocality,paperGap}]
\label{def:tower}
Write $\Ham_{d,G}$ for the condensed moduli stack of $G$-symmetric quasi-local
Hamiltonians; on a profinite probe $S$ its points are $S$-continuous families of
admissible interactions. Its \emph{uniformly gapped} substack is
\[
  \Gap_{d,G}(S)=\bigcup_{\gapbound>0}\Bigl\{\,H\in\Ham_{d,G}(S):
  \inf_{s\in S}\gapof{H_s}\ge\gapbound\,\Bigr\}.
\]
Let $\Weq$ be the class of gapped adiabatic / finite-depth quasi-local
equivalences to be inverted, and $(-)^{\st}$ stabilization by trivial
product-state ancillas. The \emph{stabilized phase $\infty$-groupoid} is
\[
  \Phase_{d,G}:=\bigl(\Gap_{d,G}[\Weq^{-1}]\bigr)^{\st},
\qquad
  \Phases_{d,G}:=\pi_0\,\Shape(\Phase_{d,G}).
\]
\end{definition}

The word \emph{uniformly} is essential: a family in which each $H_s$ is gapped,
but with $\inf_s\gapof{H_s}=0$, is excluded, and this uniformity is what
Part~III's stability theorems protect under perturbation. We take
from~\cite{paperGap} that $\Phases_{d,G}$ is well defined on the stratum where
the stability hypotheses hold with uniform constants; every statement below is
made there.

\begin{remark}[slogan]
We use the program's boxed slogan without alteration:
\begin{center}
\fbox{\parbox{0.86\linewidth}{\centering\emph{A topological phase is a component
of the stabilized condensed stack of gapped systems.}}}
\end{center}
The $\pi$-dictionary is likewise fixed: $\pi_0=$ phases; $\pi_1=$ adiabatic pumps
and phase automorphisms; $\pi_n=$ higher families and higher defects.
\end{remark}

We now fix the operation the rest of the paper is about.

\begin{definition}[stacking]
\label{def:stacking}
Let $H^{(1)},H^{(2)}$ be admissible systems on $L$ with on-site spaces
$\{h_x^{(1)}\},\{h_x^{(2)}\}$ and interactions $\Phi^{(1)},\Phi^{(2)}$. Their
\emph{stack} $H^{(1)}\boxtimes H^{(2)}$ is the system on $L$ with on-site spaces
$\{h_x^{(1)}\otimes h_x^{(2)}\}$ and interaction
\[
  (\Phi^{(1)}\boxtimes\Phi^{(2)})(X)
  =\Phi^{(1)}(X)\otimes \mathbf 1 + \mathbf 1\otimes\Phi^{(2)}(X).
\]
The \emph{trivial} system $\triv$ has one-dimensional on-site spaces and zero
interaction. Having no excited states at all, it is assigned gap $+\infty$ by
convention, so that stabilization by $\triv$ never lowers a gap
(\Cref{lem:gapstack}). Stacking of $G$-symmetric systems is $G$-symmetric under
the diagonal action.
\end{definition}

\begin{lemma}[stacking preserves the gapped substack]
\label{lem:gapstack}
If $H^{(1)},H^{(2)}$ have unique gapped ground states with gaps
$\gapbound_1,\gapbound_2>0$, then $H^{(1)}\boxtimes H^{(2)}$ has a unique gapped
ground state with
$\gapof{H^{(1)}\boxtimes H^{(2)}}=\min(\gapbound_1,\gapbound_2)$. Consequently
$\boxtimes$ restricts to a map $\Gap_{d,G}\times\Gap_{d,G}\to\Gap_{d,G}$, and if
both factors are uniformly gapped over a probe $S$ with bounds
$\gapbound_1,\gapbound_2$, the stack is uniformly gapped with bound
$\min(\gapbound_1,\gapbound_2)$.
\end{lemma}

\begin{proof}
On a finite volume $V$ the tensor-sum Hamiltonian
$H^{(1)}_V\otimes\mathbf 1+\mathbf 1\otimes H^{(2)}_V$ has spectrum
$\{\lambda+\mu\}$ with $\lambda\in\mathrm{spec}(H^{(1)}_V)$,
$\mu\in\mathrm{spec}(H^{(2)}_V)$. Its ground energy is $E^{(1)}_0+E^{(2)}_0$ with
the product ground state, unique because each factor's is; the first excited
energy is $\min\bigl(E^{(1)}_1+E^{(2)}_0,\,E^{(1)}_0+E^{(2)}_1\bigr)$, so the gap
is $\min(\gapbound_1,\gapbound_2)$. (If a factor has no excited state at all---the
trivial system $\triv$, with gap $+\infty$---the corresponding term is absent and
the minimum returns the other gap, so stacking with $\triv$ leaves the gap
unchanged.) These estimates are uniform in $V$ and,
taking infima over $s\in S$, uniform over the probe. The thermodynamic statement
follows from the infinite-volume GNS construction of~\cite{paperGap}; the minimum
of two positive uniform bounds is a positive uniform bound.
\end{proof}

That $\boxtimes$ also descends through $[\Weq^{-1}]$ and $(-)^{\st}$ is where the
symmetric-monoidal bookkeeping begins, and is the subject of the next section.

\section{Stacking as a symmetric monoidal structure}
\label{sec:monoidal}

The physics of stacking is old: put two systems side by side without coupling
them. The mathematics we need is the statement that this operation is
symmetric monoidal on the phase stack and, in particular, makes $\pi_0$ a
commutative monoid. At the level of the full $\infty$-stack the symmetric
monoidal structure is part of the program (it uses the higher stack structure of
$\Ham_{d,G}$, itself a conjecture of Part~I); the rigorous instances for
parametrized spin systems are~\cite{beaudry-etal}. At the level of $\pi_0$,
however, the monoid laws are elementary and we verify them outright.

\begin{proposition}[commutative monoid of phases]
\label{prop:monoid}
The operation induced by $\boxtimes$ makes $(\Phases_{d,G},\boxtimes,[\triv])$ a
commutative monoid: it is associative and commutative up to the equivalences in
$\Weq$, with two-sided unit the class $[\triv]$ of the trivial product state.
\end{proposition}

\begin{proof}
By \Cref{lem:gapstack}, $\boxtimes$ sends pairs of uniformly gapped families to
uniformly gapped families, so it descends to a binary operation on
$\Gap_{d,G}$. We must check it respects $\Weq$ and $(-)^{\st}$ and satisfies the
laws on classes.

\emph{Respect for $\Weq$.} If $\gamma_i:H^{(i)}\rightsquigarrow K^{(i)}$ are
gapped adiabatic / finite-depth equivalences, then
$\gamma_1\otimes\gamma_2$ is a finite-depth quasi-local equivalence
$H^{(1)}\boxtimes H^{(2)}\rightsquigarrow K^{(1)}\boxtimes K^{(2)}$: a
finite-depth circuit tensor a finite-depth circuit is finite-depth, and the
quasi-adiabatic continuation of a tensor-sum path is the tensor of the
continuations (\cite{hastings-wen-qac,bmns-automorphic}). Hence $\boxtimes$
descends to $\Gap_{d,G}[\Weq^{-1}]$ and, since stacking with $\triv$ adds trivial
ancillas, commutes with $(-)^{\st}$.

\emph{Associativity and commutativity.} On on-site spaces, $\boxtimes$ is the
tensor product of Hilbert spaces, which is associative and commutative up to the
canonical associator and the swap $\sigma_{h,h'}:h\otimes h'\to h'\otimes h$. The
tensor-sum interaction of \Cref{def:stacking} is symmetric in its two arguments
under $\sigma$, and the associator / swap are finite-depth (indeed depth-zero,
on-site) quasi-local isomorphisms, hence lie in $\Weq$. (For fermionic systems
the on-site spaces are super-vector spaces and $\sigma$ carries the Koszul sign;
this changes none of the above, as $\sigma$ is still a depth-zero on-site
isomorphism, and the symmetric-monoidal structure becomes the graded one.)
Therefore
$[H^{(1)}\boxtimes(H^{(2)}\boxtimes H^{(3)})]
 =[(H^{(1)}\boxtimes H^{(2)})\boxtimes H^{(3)}]$ and
$[H^{(1)}\boxtimes H^{(2)}]=[H^{(2)}\boxtimes H^{(1)}]$ in $\Phases_{d,G}$.

\emph{Unit.} Stacking with $\triv$ replaces each $h_x$ by $h_x\otimes\Rp\cong
h_x$ and adds nothing to the interaction, giving a depth-zero equivalence
$H\boxtimes\triv\rightsquigarrow H$ in $\Weq$. Thus $[\triv]$ is a two-sided
unit. The three laws are precisely the axioms of a commutative monoid.
\end{proof}

\begin{proposition}[invertible elements are the SRE phases]
\label{prop:sre}
An element $[H]\in\Phases_{d,G}$ is invertible for $\boxtimes$ if and only if
there is a system $\bar H$ with $[H]\boxtimes[\bar H]=[\triv]$, i.e.\ $H$ is
\emph{short-range entangled}: stacked with a partner it becomes trivial after
stabilization. Write $\inv{d,G}\subseteq\Phase_{d,G}$ for the full sub-$\infty$-groupoid
on the invertible objects and $\Phases_{d,G}^{\times}$ for its $\pi_0$.
\end{proposition}

\begin{proof}
This is the definition of an invertible element in a monoid, transported through
$\pi_0$. That the invertible objects form a full symmetric-monoidal
sub-$\infty$-groupoid is formal: invertibility is preserved by $\boxtimes$
(a tensor of invertibles is invertible with inverse the tensor of inverses) and
detected on $\pi_0$; the equivalences in $\Weq$ preserve it.
\end{proof}

The candidate partner $\bar H$ is, physically, the orientation-reversed or
complex-conjugate system; that every SRE phase actually admits such an inverse is
the content of the invertibility of the low-energy TQFT
\cite{freed-sre,freed-hopkins} and is used, not proved, here. The distinction
between the full monoid $\Phases_{d,G}$ and its invertible submonoid
$\Phases_{d,G}^{\times}$ is the distinction between ``all gapped phases''
(including intrinsic topological order, which has no inverse) and the SRE phases
that an invertible EFT can see. The spectrum of \Cref{sec:spectrum} is built from
the latter; the group completion of \Cref{sec:groupcompletion} is what the former
needs.

\section{Group completion and the stabilization element}
\label{sec:groupcompletion}

A commutative monoid is not a group: some phases have no inverse, and even among
those that ``cancel'' the cancellation may require adding a common summand. The
universal way to force inverses is group completion, and its explicit form is
exactly where the physics of stabilization by ancillas becomes visible. Nothing
in this section is new mathematics: it is the Grothendieck construction, but
the reading of the stabilization element is the point.

\begin{definition}[group completion]
\label{def:groth}
Let $(M,+,0)$ be a commutative monoid. Its \emph{group completion} (Grothendieck
group) is $\Groth(M):=(M\times M)/{\sim}$, where $(a,b)\sim(a',b')$ iff there is
$e\in M$ with
\[
  a+b'+e \;=\; a'+b+e .
\]
Addition is $[(a,b)]+[(c,d)]=[(a+c,b+d)]$, and
$\gamma_M:M\to\Groth(M)$, $\gamma_M(m)=[(m,0)]$.
\end{definition}

\begin{theorem}[\textbf{IV-A}: universal property of group completion]
\label{thm:groth}
For any commutative monoid $M$:
\begin{enumerate}
\item[(i)] $\sim$ is a congruence and $\Groth(M)$ is an abelian group, with
  identity $[(0,0)]$ and $-[(a,b)]=[(b,a)]$;
\item[(ii)] $\gamma_M$ is a monoid homomorphism;
\item[(iii)] for every abelian group $A$ and monoid homomorphism $f:M\to A$
  there is a unique group homomorphism $\bar f:\Groth(M)\to A$ with
  $\bar f\circ\gamma_M=f$.
\end{enumerate}
The pair $(\Groth(M),\gamma_M)$ is thereby determined up to unique isomorphism.
\end{theorem}

\begin{proof}
\emph{(i)} Reflexivity and symmetry are immediate. For transitivity, suppose
$(a,b)\sim(c,d)$ via $e$ and $(c,d)\sim(g,h)$ via $e'$, so
$a+d+e=c+b+e$ and $c+h+e'=g+d+e'$. Put $e''=d+e+e'$. Then
\begin{align*}
a+h+e'' &= (a+d+e)+h+e' = (c+b+e)+h+e' = b+e+(c+h+e')\\
        &= b+e+(g+d+e') = g+b+e'',
\end{align*}
so $(a,b)\sim(g,h)$. If $(a,b)\sim(a',b')$ via $e$ then adding $c+d$ shows
$(a+c,b+d)\sim(a'+c,b'+d)$ via the same $e$, so addition is well defined; it
inherits associativity, commutativity, and the identity $[(0,0)]$ from $M$.
Finally $[(a,b)]+[(b,a)]=[(a+b,a+b)]=[(0,0)]$ because $(a+b,a+b)\sim(0,0)$ via
$e=0$. Hence $\Groth(M)$ is an abelian group.

\emph{(ii)} $\gamma_M(m)+\gamma_M(n)=[(m+n,0)]=\gamma_M(m+n)$ and
$\gamma_M(0)=[(0,0)]$.

\emph{(iii)} Set $\bar f([(a,b)]):=f(a)-f(b)\in A$. This is well defined: if
$(a,b)\sim(a',b')$ via $e$ then $f(a)+f(b')+f(e)=f(a')+f(b)+f(e)$ in $A$, and $A$
is a group, so $f(a)-f(b)=f(a')-f(b')$. It is a homomorphism by construction and
$\bar f(\gamma_M(m))=f(m)-f(0)=f(m)$. For uniqueness, note
$[(a,b)]=\gamma_M(a)-\gamma_M(b)$ in $\Groth(M)$; any $g$ with
$g\circ\gamma_M=f$ therefore satisfies
$g([(a,b)])=f(a)-f(b)=\bar f([(a,b)])$.
\end{proof}

The construction is functorial: a monoid homomorphism $\varphi:M\to N$ induces
$\Groth(\varphi):\Groth(M)\to\Groth(N)$, and $\Groth$ is left adjoint to the
forgetful functor from abelian groups to commutative monoids. We will use this
adjunction when discussing whether solidification commutes with group completion
(Conjecture IV-4).

Now the element $e$. Its presence is the whole difference between a group and a
monoid, and it is the algebraic shadow of ancilla stabilization.

\begin{proposition}[the stabilization element]
\label{prop:stable}
For $a,b\in M$, $\gamma_M(a)=\gamma_M(b)$ in $\Groth(M)$ if and only if there is
$e\in M$ with $a+e=b+e$. Consequently $\gamma_M$ is injective iff $M$ is
cancellative, and in general the largest quotient of $M$ on which $\gamma_M$ is
injective is $M^{\st}:=M/{\approx}$, where $a\approx b$ iff $a+e=b+e$ for some
$e$. The monoid $M^{\st}$ is cancellative and
$\Groth(M)\cong\Groth(M^{\st})$.
\end{proposition}

\begin{proof}
$\gamma_M(a)=\gamma_M(b)$ means $(a,0)\sim(b,0)$, i.e.\ $a+0+e=b+0+e$ for some
$e$; that is $a+e=b+e$. Injectivity of $\gamma_M$ is then the implication
$a+e=b+e\Rightarrow a=b$, i.e.\ cancellativity. The relation $\approx$ is a
congruence (if $a+e=b+e$ then $a+c+e=b+c+e$ for every $c$), and $M^{\st}$ is
cancellative because, for an \emph{arbitrary} $[c]\in M^{\st}$,
$[a]+[c]=[b]+[c]$ in $M^{\st}$ means $a+c+e'=b+c+e'$ for some $e'$, hence
$a+e''=b+e''$ with $e''=c+e'$, i.e.\ $a\approx b$. Since $\approx$ is exactly the
kernel pair of $\gamma_M$, the map
$M^{\st}\to\Groth(M)$ is injective and induces
$\Groth(M^{\st})\xrightarrow{\ \cong\ }\Groth(M)$ by the universal property.
\end{proof}

\begin{remark}[why $(-)^{\st}$ is in \Cref{def:tower}, and what it is not]
\label{rem:stmeaning}
\Cref{prop:stable} is the reason the stabilization $(-)^{\st}$ appears in the
definition of $\Phase_{d,G}$: two gapped systems are declared the same phase when
they agree \emph{after adding trivial product-state ancillas}. This is a genuine
operation, not stacking with the unit $\triv$. A product state on
higher-dimensional on-site spaces is not $\Weq$-equivalent to $\triv$ (finite-depth
quasi-local circuits preserve on-site Hilbert-space dimension, so they cannot
shrink an ancilla back to a point), and inverting ``$\boxtimes$ a trivial product
state'' therefore enlarges the equivalence in a way stacking with $\triv$ never
could. Write $\approx_{\mathrm{prod}}$ for the resulting relation
($a\approx_{\mathrm{prod}}b$ iff $a+e=b+e$ for some \emph{product state} $e$); the
physically stabilized phase monoid is $M/{\approx_{\mathrm{prod}}}$.

Two cautions keep this honest. First, $\approx_{\mathrm{prod}}$ is in general
\emph{finer} than the full stable equality $\approx$ of \Cref{prop:stable}, which
allows an \emph{arbitrary} $e\in M$: when intrinsic topological order is present,
$a+m=b+m$ for a non-invertible $m$ need not entail $a+e=b+e$ for any product
state $e$. So $M/{\approx_{\mathrm{prod}}}$ need not be cancellative, and the
group completion $\Groth$ can collapse it \emph{further}; the natural map
$M/{\approx_{\mathrm{prod}}}\to\Groth(\Phases_{d,G})$ is then not injective. This
is the ``collapse'' discussed in \Cref{sec:discussion}, and it is why we do
\emph{not} claim the full phase monoid embeds in its group completion. Second, on
the \emph{invertible} submonoid the two relations coincide: if $a$ is invertible
with partner $\bar a$, then $a\boxtimes\bar a$ is stably a product state, so any
witness $e$ of $a+e=b+e$ can be traded for a product state after stacking with
enough copies of $\bar a$. A grouplike monoid is already cancellative and equals
its own group completion, so
$\Phases_{d,G}^{\times}\cong\Groth(\Phases_{d,G}^{\times})$ is an abelian group
with no further completion needed, which is exactly why the spectrum of
\Cref{sec:spectrum} is built from the invertible sector, the one place where
stabilization and group completion agree.
\end{remark}

\begin{example}[a non-cancellative toy]
\label{ex:noncanc}
Let $M=\langle a,t\mid t+a=t\rangle$ be the commutative monoid on generators
$a,t$ with the single relation $t+a=t$ (a generator $t$ that absorbs $a$). Here
$a\ne 0$ in $M$, yet $\gamma_M(a)=\gamma_M(0)$ because $a+t=0+t$: the
stabilization element is $e=t$. So $a$ becomes invisible in $\Groth(M)$---the
algebra behind ``a phase nontrivial on the nose but trivial once enough ancillas
are added.'' This does \emph{not} make the completion vanish: no relation forces
$t$ to be absorbed---the only relation is $t+a=t$, and $nt+e=e$ has no solution
in $M$ for $n\ge1$---so the multiples of $t$ are cancellative among themselves,
and
\[
  \Groth(M)\;\cong\;\Zp,\qquad\text{generated by }[t],\quad\text{with }[a]=0 .
\]
The example is non-cancellativity in its purest form---one generator collapses
under stabilization while another survives. The accompanying code
(\Cref{sec:haskell}) computes $\Groth(M)$ for this $M$ and confirms both facts:
$[a]=0$, while the multiples $[nt]$ ($n\ge1$) are pairwise distinct.
\end{example}

\section{The invertible condensed phase spectrum}
\label{sec:spectrum}

We now assemble the top floor of \eqref{eq:tower}. The mechanism is the standard
one that turns a grouplike commutative monoid object in spaces into a connective
spectrum; what is specific to our setting is the identification of the input with
the invertible phase groupoid and the (conjectural) condensed refinement.

\subsection{From the invertible sector to a connective spectrum}

By \Cref{prop:sre} the invertible sector $\inv{d,G}$ is a symmetric-monoidal
sub-$\infty$-groupoid of $\Phase_{d,G}$ in which every object is
$\boxtimes$-invertible; that is, a \emph{Picard $\infty$-groupoid}, equivalently
a grouplike $E_\infty$-space.

\begin{theorem}[recognition, conditional]
\label{thm:recognition}
Suppose $\inv{d,G}$ is a Picard $\infty$-groupoid (a grouplike $E_\infty$-space).
Then there is a connective spectrum $\IPcond_{d,G}$, unique up to equivalence,
with $\Omega^\infty\IPcond_{d,G}\simeq\inv{d,G}$, and
\[
  \pi_0\IPcond_{d,G}\cong\Phases_{d,G}^{\times},\qquad
  \pi_n\IPcond_{d,G}\cong\pi_n\bigl(\inv{d,G}\bigr)\ (n\ge 1).
\]
In particular $\pi_1$ records adiabatic pumps / phase automorphisms and $\pi_n$
higher families, per the $\pi$-dictionary.
\end{theorem}

\begin{proof}
This is the recognition principle for grouplike $E_\infty$-spaces: the
$\infty$-category of connective spectra is equivalent to that of grouplike
$E_\infty$-spaces via $\Omega^\infty$, with inverse the connective-spectrum
functor. The hypothesis is that $\inv{d,G}$ is such an object. We do not reprove
the recognition principle; in the precise setting of quantum lattice systems the
passage from the invertible / grouplike space of gapped systems to a spectrum is
carried out rigorously in~\cite{beaudry-etal} (parametrized foundations) and
\cite{kubota-omega-spectrum} (an explicit $\Omega$-spectrum). We record the
statement conditionally because the hypothesis---that $\Phase_{d,G}$ is a
symmetric-monoidal $\infty$-groupoid with the invertible sector grouplike---rests
on the higher-stack structure of $\Ham_{d,G}$, which is Conjecture~I-1 of
\cite{paperLocality}.
\end{proof}

\begin{definition}[invertible condensed phase spectrum]
\label{def:ipcond}
$\IPcond_{d,G}$ is the connective spectrum of \Cref{thm:recognition}, the
\emph{invertible condensed phase spectrum}. Its defining property is
$\Omega^\infty\IPcond_{d,G}\simeq\inv{d,G}$.
\end{definition}

The name contains a promissory note: as constructed, $\IPcond_{d,G}$ is a
spectrum of \emph{spaces} (the shape of the invertible condensed groupoid), and
whether it genuinely lifts to a spectrum object internal to condensed / solid
mathematics is Conjecture IV-1 below. We keep the notation $\IPcond$ throughout
and flag the conjectural status wherever it matters.

\subsection{A rigorous carrier for the homotopy groups}

The value of \Cref{thm:recognition} is limited by its hypothesis. What rescues it
from vacuity is that its output has an independently constructed, fully rigorous
model.

\begin{theorem}[\textbf{IV-B}: Kubota's $\Omega$-spectrum,~\cite{kubota-omega-spectrum}]
\label{thm:kubota}
There is an $\Omega$-spectrum $\mathit{IP}^\ast$, built from the operator-algebraic
formulation of invertible gapped quantum spin systems, whose homotopy groups are
the groups of invertible gapped spin systems in each dimension (with variants for
crystallographic symmetry). It realizes Kitaev's proposal
\cite{kitaev-periodic-table} that invertible phases are the homotopy groups of a
spectrum.
\end{theorem}

This is a cited theorem, not ours. Its relationship to \Cref{def:ipcond} is the
subject of the next two conjectures: $\mathit{IP}^\ast$ is the smooth / operator-algebraic
avatar of the spaces underlying $\IPcond_{d,G}$, and the content of Conjecture
IV-1 is that $\IPcond$ is its condensed refinement.

\begin{conjecture}[\textbf{IV-1}: condensed / solid refinement]
\label{conj:iv1}
$\IPcond_{d,G}$ is a connective spectrum object internal to solid modules
(equivalently, a condensed connective spectrum) whose underlying spectrum of
spaces---its shape, obtained by forgetting the solid structure---is a connective
cover of Kubota's $\mathit{IP}^\ast$~\cite{kubota-omega-spectrum}, and whose
homotopy solid modules refine the homotopy groups of $\mathit{IP}^\ast$. The
refinement is the one carried
by the solid $K$-theory of the condensed observable algebra of Part~II
\cite{paperPositivity}, via Aoki's identification
$\Kop(A)\simeq\Solid(\Kalg(\cond{A}))$~\cite{aoki-solidification}.
\end{conjecture}

The reason to want IV-1, rather than to be content with Kubota's smooth spectrum,
is uniformity across the program: the profinite-disorder hulls $\dis=Q^{\Zp^d}$,
the continuous families over condensed probes, and the analytic completions all
live naturally in condensed / solid mathematics
\cite{clausen-scholze-condensed,clausen-scholze-analytic,clausen-scholze-analytic-stacks},
and a spectrum internal to that world is what makes level (L2) of
\Cref{subsec:levels} even statable. It is, however, a genuine conjecture: no
condensed enhancement of $\mathit{IP}^\ast$ exists in the literature.

\section[Renormalization and the effective-field-theory passage]{Renormalization and the\\effective-field-theory passage}
\label{sec:rg}

Between a lattice system and its EFT sits renormalization: coarse-grain, discard
short-distance data, iterate, take a limit. We formulate this as a filtered /
pro-system and identify its analytic limit with solidification. Both statements
are conjectural; the point is to say precisely \emph{what} would have to be true.
The two conjectures of this section carry the stable series identifiers IV-5 and
IV-4: like every IV-$n$ label in this paper they are fixed across the six-part
series and referenced by the same name elsewhere, so they are not ordered by
first appearance here (IV-2, the central comparison conjecture, follows in
\Cref{sec:comparison}).

\subsection{Coarse-graining as a filtered system}

\begin{definition}[block-spin coarse-graining]
\label{def:coarse}
A \emph{coarse-graining} of scale $\ell$ is a map
$R_\ell:\Ham_{d,G}\to\Ham_{d,G}$ that partitions $L$ into blocks of diameter
$\ell$, replaces the on-site space of a block by a chosen subspace (a truncation
or isometry $V_\ell$), and pushes the interaction forward, $\Phi\mapsto
V_\ell^\ast\Phi V_\ell$, retaining the induced quasi-local structure. Composition
of scales gives a filtered system
$\cdots\to R_{\ell_2}\to R_{\ell_1}\to\mathrm{id}$ indexed by the poset of
scales.
\end{definition}

The finite-resolution truncations of a profinite disorder hull
$\dis=\varprojlim\dis_i$ (\cite{paperPositivity}, \S1.7) are the disorder-theoretic
instance of exactly this poset: coarse-graining in space and coarsening the
configuration data are the same kind of pro-operation. This is why the EFT limit
and the descent along $\dis$ share a formalism.

\begin{conjecture}[\textbf{IV-5}: renormalization functor]
\label{conj:iv5}
The coarse-grainings $\{R_\ell\}$ assemble into a filtered system of endofunctors
of $\Gap_{d,G}$ compatible with $\boxtimes$ (a lax symmetric-monoidal
pro-endofunctor), whose limit / colimit
\[
  R_\infty:=\colim_\ell R_\ell
\]
exists on the SRE stratum and computes the effective field theory: for an SRE
family $H$, $R_\infty H$ is (the lattice model of) the fixed-point theory whose
deformation class is the EFT invariant. The induced map on phases
$\Phases_{d,G}\to\Phases_{d,G}$ is idempotent with image the SRE sub-monoid.
\end{conjecture}

Coarse-graining is delicate for the same reason the gap is: the undecidability of
the spectral gap~\cite{cpw-undecidable} forbids any general algorithm certifying
that a coarse-graining flows to a gapped fixed point, so IV-5 is asserted on the
SRE stratum only, and even there its content is the existence of the limit, not a
procedure to compute it. This mirrors the conditional stance Part~III is forced
into.

\subsection{Solidification as the analytic completion}

At the level of invariants there is a candidate for $R_\infty$ that is already
rigorous as a functor, if not yet known to model coarse-graining: Aoki's
solidification. Part~II installs the identification
$\Kop(A)\simeq\Solid(\Kalg(\cond{A}))$~\cite{aoki-solidification}: the operator
$K$-theory that carries topological-insulator invariants is the
\emph{solidification} of the algebraic $K$-theory of the condensed observable
algebra. Solidification is a completion: it inverts the analytic / topological
data that algebraic $K$-theory does not see, and this is exactly the role
``passing to the EFT'' plays for invariants: discard the lattice-scale algebraic
information, keep the topological class.

\begin{conjecture}[\textbf{IV-4}: solidification commutes with group completion]
\label{conj:iv4}
Let $M_{d,G}$ be the condensed commutative monoid of stacking classes and
$\Groth$ its group completion (\Cref{thm:groth}) computed internally to condensed
abelian groups. Then solidification commutes with group completion,
\[
  \Solid\bigl(\Groth(M_{d,G})\bigr)\;\simeq\;\Groth\bigl(\Solid(M_{d,G})\bigr),
\]
naturally in $(d,G)$. Equivalently, the analytic completion (solidification) and
the algebraic completion (group completion / adding formal inverses) are
compatible, so that the invertible solid spectrum may be built in either order.
\end{conjecture}

Conjecture IV-4 is the compatibility that makes ``$\IPcond$ built from lattice
data'' and ``$\IPcond$ built from solid $K$-theory'' agree. It is plausible on
formal grounds---$\Groth$ is a left adjoint (\Cref{thm:groth}) and solidification
is a localization / left adjoint~\cite{clausen-scholze-analytic}, and left
adjoints compose---but the two adjunctions live in different categories (condensed
monoids versus solid modules) and the interchange is not automatic. We state it as
a conjecture rather than dress a non-theorem as a corollary.

\subsection{What ``equivalence'' can honestly mean here}

We can now say precisely what the EFT passage delivers at each of the three
levels of \Cref{subsec:levels}.
\begin{itemize}
\item At \textbf{(L0)}, the EFT passage is the idempotent
  $\Phases_{d,G}\to\Phases_{d,G}^{\times}$ of Conjecture IV-5 followed by the
  microscopic-to-field-theoretic identification of Conjecture IV-2; ``lattice
  $=$ EFT'' is a statement about which \emph{set} of phases one gets.
\item At \textbf{(L1)}, it is the equivalence of spectra $\IPcond_{d,G}\simeq
  \IFH$ after solidification, requiring Conjectures IV-1 and IV-4 to even phrase
  the completion.
\item At \textbf{(L2)}, it is naturality over condensed probes, requiring the
  condensed refinement of IV-1 in an essential way.
\end{itemize}
None of these is proved. The honest content of this section is the
identification of the missing pieces and the fact that the analytic completion
step is not mysterious: it is solidification, a functor we already have.

\section{The comparison conjecture}
\label{sec:comparison}

We reach the organizing statement. On one side is $\IPcond_{d,G}$, built from
lattice systems; on the other is the Freed--Hopkins classification of invertible
field theories, built from bordism and Anderson duality. The comparison
conjecture is that they agree.

\subsection{The field-theory target}

\begin{definition}[Freed--Hopkins target,~\cite{freed-hopkins}]
\label{def:fh}
Fix a symmetry type $H_d$ (a stable tangential structure, e.g.\ $\mathrm{Spin}$,
$\mathrm{Spin}\times G$, $\mathrm{Pin}^\pm$) with Madsen--Tillmann spectrum
$\MTH$. Reflection-positive invertible $d$-dimensional field theories with
symmetry type $H_d$ are classified by the abelian group
\[
  \bigl[\MTH,\ \Sigma^{d+1}\IZ\bigr],
\]
homotopy classes of spectrum maps into a shift of the Anderson dual $\IZ$ of the
sphere. Assembling over $d$ gives a spectrum $\IFH$ whose homotopy groups are
these deformation-class groups; its torsion part is the ``beyond group
cohomology'' content~\cite{kapustin-cobordism} and its free part records the
integer invariants (Hall conductances, chiral central charges).
\end{definition}

The passage from a short-range-entangled state to such a theory is Freed's
theorem that an SRE system defines an invertible field theory
\cite{freed-sre}, and the generalized-cohomology / $\Omega$-spectrum viewpoint of
\cite{gaiotto-jf,xiong-minimalist} is the statement that these deformation
classes are themselves the homotopy of a spectrum---the field-theoretic sibling
of \Cref{thm:recognition}.

\subsection{The comparison map and the conjecture}

A lattice invertible phase has a low-energy theory; taking its deformation class
defines a comparison map. We isolate it as a hypothesis and then conjecture it is
an equivalence.

\begin{definition}[comparison map, conditional]
\label{def:comp}
Assume the recognition hypothesis of \Cref{thm:recognition}---that
$\Phase_{d,G}$ is a symmetric-monoidal $\infty$-groupoid, equivalently
Conjecture~I-1 of~\cite{paperLocality}---so that the invertible sector assembles
into the spectrum $\IPcond_{d,G}$; and assume Freed's short-range-entanglement
assignment extends naturally in $d$ and in the symmetry data, a naturality we do
not construct here. Under these hypotheses one obtains a \emph{comparison map} of
spectra
\[
  c_{d,G}:\ \IPcond_{d,G}\ \longrightarrow\ \IFH,
\]
induced on invertible sectors by sending an SRE lattice family to the
deformation class of its low-energy invertible field theory (Freed's assignment
\cite{freed-sre}). Absent the spectrum-level hypotheses, only the induced map on
components $\pi_0(c_{d,G}):\Phases_{d,G}^{\times}\to\pi_0\IFH$ is unconditional;
it is this $\pi_0$ map that the (L0) evidence of \Cref{sec:evidence} constrains.
\end{definition}

\begin{conjecture}[\textbf{IV-2}: lattice--EFT comparison equivalence]
\label{conj:iv2}
Under short-range-entanglement (EFT) hypotheses, the comparison map $c_{d,G}$ is
an equivalence after an appropriate completion. Explicitly, at the three levels
of \Cref{subsec:levels}:
\begin{enumerate}
\item[\textbf{(L0)}] $\pi_0(c_{d,G}):\Phases_{d,G}^{\times}\xrightarrow{\ \cong\ }
  \pi_0\IFH$ is an isomorphism of abelian groups: the microscopic invariant
  equals the deformation class of the associated invertible TQFT;
\item[\textbf{(L1)}] $c_{d,G}$ is an equivalence of connective spectra after
  solidification / on the relevant completion;
\item[\textbf{(L2)}] $c_{d,G}$ upgrades to a natural equivalence of condensed
  spectra over profinite probes, compatible with disorder hulls $\dis$.
\end{enumerate}
In particular the SSH chain matches its Dirac EFT (\Cref{ex:ssh}).
\end{conjecture}

This is the paper's central claim and it is \emph{open}. The remainder of the
paper is evidence and worked cases, never a proof. It is worth being explicit
about the two ways IV-2 could fail even for SRE systems: the comparison could be
injective but not surjective (some field theory not realized by any lattice
model---this is precisely the realizability question handed to Part~V
\cite{paperRealizability}, where the Kapustin--Fidkowski obstruction
\cite{kapustin-fidkowski} shows commuting-projector models cannot realize chiral
classes), or surjective but not injective (distinct lattice phases with the same
EFT---expected to be prevented by stabilization but not proven in general).

\subsection{Structure of the comparison}

The comparison sits in a square that summarizes the paper:
\[
\begin{tikzcd}[column sep=large, row sep=large]
\{\text{SRE lattice families}\}/\Weq \arrow[r, "\text{stabilize}"] \arrow[d, "\text{low-energy}"']
  & \Phases_{d,G}^{\times}\arrow[d, "\pi_0(c_{d,G})"] \\
\{\text{invertible TQFTs}\}/{\simeq} \arrow[r, "\text{deform.\ class}"']
  & \pi_0\IFH
\end{tikzcd}
\]
The left vertical is Freed's SRE-to-TQFT assignment; the right vertical is the
map IV-2 conjectures to be an isomorphism; the horizontals are the two
classifications. Commutativity of the square is definitional; the conjecture is
that the right vertical is an isomorphism, whence---given the outer maps---the
two notions of ``phase'' coincide.

\section{Evidence}
\label{sec:evidence}

Three bodies of rigorous work constrain Conjecture IV-2. In one regime it is a
theorem; in another it is realized by classical $K$-theory; in a third it has a
rigorous homotopical target. We present each honestly, marking exactly how far it
reaches.

\subsection{\texorpdfstring{$d=1$}{d=1}: completeness makes (L0) a theorem}

\begin{theorem}[$d=1$ comparison, from~\cite{ogata-spt-chains,ogata-classification-review}]
\label{thm:d1}
For one-dimensional quantum spin chains with on-site finite symmetry $G$, the
operator-algebraic index built from the split property is a \emph{complete}
invariant of SPT phases, valued in $H^2(G,\Tt)$. Since $H^2(G,\Tt)$ is exactly
the field-theoretic (group-cohomology) classification of $1{+}1$-dimensional
$G$-SPTs, the level-(L0) comparison $\pi_0(c_{1,G})$ is a bijection. In this
sense Conjecture IV-2 holds at (L0) in $d=1$.
\end{theorem}

\begin{proof}[Discussion of proof]
The completeness of the operator-algebraic index for $d=1$ on-site finite
symmetry is Ogata's theorem~\cite{ogata-spt-chains} (surveyed
in~\cite{ogata-classification-review}); its value group
$H^2(G,\Tt)$ coincides with the group-cohomology classification of
$(1{+}1)$D bosonic SPTs~\cite{cglw-cohomology}, which is the deformation-class
group $\pi_0\IFH$ in this case. The identification of the two---that the index of
a chain equals the class of its EFT---is the matching used in the
operator-algebraic literature; we cite rather than reprove it. Note the
hypotheses are exactly the EFT hypotheses of IV-2 specialized to $d=1$: on-site
finite symmetry, unique gapped ground state.
\end{proof}

This is the strongest existing anchor: a regime where ``lattice $=$ EFT'' is a
theorem at the level of phase sets. It does not touch (L1) or (L2)---the split-property
index is a $\pi_0$ statement---which is exactly why we separated the levels. In
$d=2$ Ogata's $H^3(G,\Tt)$-valued index~\cite{ogata-h3-index,ogata-spt-chains}
provides a well-defined bulk invariant that matches the expected group-cohomology
label; but its \emph{completeness}---that it separates all $2$D SPT phases with
on-site finite symmetry, as the split-property index does in $d=1$---is not
established. So the (L0) comparison in $d=2$ is at present partial: a realized
invariant, not a settled bijection.

\subsection{Free fermions: the tenfold way realizes the comparison}

For free-fermion systems the comparison is classical and complete at (L0) in
every dimension, because both sides are computed by the same $K$-theory. This is
the Kitaev periodic table~\cite{kitaev-periodic-table}, whose mathematical
content is $K$-theoretic~\cite{thiang-ktheory,prodan-sb}.

\begin{definition}[the tenfold way]
\label{def:tenfold}
The ten Altland--Zirnbauer symmetry classes~\cite{altland-zirnbauer} split into
two complex classes ($A,\mathit{AIII}$) and eight real classes
($\mathit{AI},\mathit{BDI},D,\mathit{DIII},\mathit{AII},\mathit{CII},C,\mathit{CI}$),
indexed by a symmetry label $s\in\Zp/2$ (complex) or $s\in\Zp/8$ (real). In
spatial dimension $d$, the group of strong topological invariants of gapped
free-fermion Hamiltonians in the class is
\[
  \mathcal G_{\mathbb C}(s,d)=d_{\mathbb C}\bigl((s-d)\bmod 2\bigr),
  \qquad
  \mathcal G_{\mathbb R}(s,d)=d_{\mathbb R}\bigl((s-d)\bmod 8\bigr),
\]
where the complex and real Bott sequences are
\[
  d_{\mathbb C}:\ (\Zp,\,0),\qquad
  d_{\mathbb R}:\ (\Zp,\,\Zp/2,\,\Zp/2,\,0,\,2\Zp,\,0,\,0,\,0),
\]
read cyclically at indices $0,1,\dots$. These are the homotopy groups of the
$K$-theory spectra $KU$ and $KO$ up to the standard $2\Zp$ labelling of the
degree-$4$ real generator.
\end{definition}

\begin{proposition}[Bott periodicity and the dimension shift]
\label{prop:bott}
The tenfold-way invariant groups of \Cref{def:tenfold} are Bott periodic:
$\mathcal G_{\mathbb C}(s,d)=\mathcal G_{\mathbb C}(s,d+2)$ and
$\mathcal G_{\mathbb R}(s,d)=\mathcal G_{\mathbb R}(s,d+8)$. They are moreover
invariant under the simultaneous shift $(s,d)\mapsto(s+1,d+1)$, so the whole table
is determined by the single antidiagonal $s-d$; equivalently, raising the
symmetry label and the dimension together leaves the phase group unchanged. These
periodicities are the free-fermion incarnation of the spectrum structure of
\Cref{thm:recognition}, with the dimension shift $d\mapsto d+1$ acting as the
degree shift (multiplication by the Bott generator).
\end{proposition}

\begin{proof}
Both groups are defined by reduction of $s-d$ modulo the period ($2$ or $8$), so
periodicity in $d$ and the antidiagonal invariance $(s,d)\mapsto(s+1,d+1)$ are
immediate from
$\mathcal G(s,d)=d_\bullet\bigl((s-d)\bmod p\bigr)$. That the resulting groups
are the correct classifying groups of the ten classes---and that the
$d\mapsto d+1$ shift is Bott periodicity of $KO$/$KU$---is the content of
\cite{kitaev-periodic-table,thiang-ktheory}; we use their identification. The
periodicities are verified independently by machine in \Cref{sec:haskell}.
\end{proof}

\begin{table}[t]
\centering
\renewcommand{\arraystretch}{1.25}
\begin{tabular}{lcccccccc}
\toprule
class ($s$) & $d{=}0$ & $1$ & $2$ & $3$ & $4$ & $5$ & $6$ & $7$\\
\midrule
$\mathit{AI}$ (0)   & $\Zp$   & $0$    & $0$    & $0$    & $2\Zp$ & $0$    & $\Zp/2$& $\Zp/2$\\
$\mathit{BDI}$ (1)  & $\Zp/2$ & $\Zp$  & $0$    & $0$    & $0$    & $2\Zp$ & $0$    & $\Zp/2$\\
$D$ (2)             & $\Zp/2$ & $\Zp/2$& $\Zp$  & $0$    & $0$    & $0$    & $2\Zp$ & $0$\\
$\mathit{DIII}$ (3) & $0$     & $\Zp/2$& $\Zp/2$& $\Zp$  & $0$    & $0$    & $0$    & $2\Zp$\\
$\mathit{AII}$ (4)  & $2\Zp$  & $0$    & $\Zp/2$& $\Zp/2$& $\Zp$  & $0$    & $0$    & $0$\\
$\mathit{CII}$ (5)  & $0$     & $2\Zp$ & $0$    & $\Zp/2$& $\Zp/2$& $\Zp$  & $0$    & $0$\\
$C$ (6)             & $0$     & $0$    & $2\Zp$ & $0$    & $\Zp/2$& $\Zp/2$& $\Zp$  & $0$\\
$\mathit{CI}$ (7)   & $0$     & $0$    & $0$    & $2\Zp$ & $0$    & $\Zp/2$& $\Zp/2$& $\Zp$\\
\midrule
$A$ (0)             & $\Zp$   & $0$    & $\Zp$  & $0$    & $\Zp$  & $0$    & $\Zp$  & $0$\\
$\mathit{AIII}$ (1) & $0$     & $\Zp$  & $0$    & $\Zp$  & $0$    & $\Zp$  & $0$    & $\Zp$\\
\bottomrule
\end{tabular}
\caption{The tenfold-way table of \Cref{def:tenfold}, entry
$\mathcal G(s,d)$. Real classes (top eight) are $8$-periodic along each row and
constant along antidiagonals $s-d=\text{const}$; complex classes (bottom two) are
$2$-periodic. Landmark entries: $D$ at $d{=}1$ (Kitaev chain, $\Zp/2$); $D$ at
$d{=}2$ ($p{+}ip$, $\Zp$); $\mathit{AII}$ at $d{=}2,3$ (quantum spin Hall / 3D
TI, $\Zp/2$); $A$ at $d{=}2$ (integer quantum Hall, $\Zp$); $\mathit{AIII}$ at
$d{=}1$ (SSH, $\Zp$).}
\label{tab:tenfold}
\end{table}

\Cref{tab:tenfold} lists the invariants. Every landmark free-fermion phase sits
at its expected entry, and the table is the concrete low-dimensional
consistency check the program asks for: in each of these boxes the microscopic
$K$-theory invariant equals the invariant of the massive-Dirac EFT, so
Conjecture IV-2 holds at (L0) for free fermions throughout. The free-fermion case
also gives the cleanest picture of the dimension-reduction functoriality: the
antidiagonal invariance of \Cref{prop:bott} is the statement that adding a
symmetry and adding a dimension cancel, which is dimensional reduction in the
sense of~\cite{prodan-sb}.

\subsection{\texorpdfstring{$\Omega$}{Omega}-spectra: a rigorous target for (L1)}

Finally, the homotopical anchors. Kubota's $\Omega$-spectrum
(\Cref{thm:kubota}) and the parametrized foundations of~\cite{beaudry-etal}
provide, for the first time, a rigorous \emph{target} at level (L1): a spectrum
whose homotopy groups are the invertible gapped systems, against which
$c_{d,G}$ can be compared as a map of spectra. They do not prove IV-2 at (L1)
(that would require identifying Kubota's spectrum with the Freed--Hopkins
target, which is itself open), but they turn (L1) from a slogan into a comparison
of two specified spectra. This is the sense in which the program is ``rigorous at
the module level, conjectural at the global level'': the objects are real, the
identification is conjectural.

\section{Lattice invariants, the SSH example, and transitions}
\label{sec:invariants}

The comparison map $c_{d,G}$ is abstract; the invariants that instantiate it are
concrete. We record a rigorous lattice invariant, run the SSH example through the
whole apparatus, and connect transitions to the spectrum's boundary map.

\subsection{A rigorous lattice invariant}

\begin{theorem}[\textbf{IV-C}: Hall conductance as a phase invariant,~\cite{kapustin-sopenko-hall,kapustin-sopenko-noether}]
\label{thm:hall}
For a two-dimensional short-range-entangled lattice state, the Hall conductance
is locally computable from the state, is an integer multiple of $e^2/h$, and is
constant on gapped phases; it is therefore a well-defined homomorphism
$\Phases_{2,G}^{\times}\to\Zp$. For families over a base, the higher Berry class
of~\cite{kapustin-sopenko-noether} generalizes it and unifies Hall conductance
with the Thouless charge pump~\cite{thouless-pump}.
\end{theorem}

This is exactly a component of the comparison map: $\Phases_{2,G}^{\times}\to\Zp
=\pi_0\IFH$ in class $A$ (integer quantum Hall), the free part of
\Cref{def:fh}, computed microscopically from the lattice state. That it is
locally computable and integer-quantized is the rigorous content behind ``the
lattice invariant equals the EFT invariant'' in this box, and it is the
generalized-cohomology invariant $\nu\in E^q(U_f)$ of the program's transition
calculus. The Kapustin--Sopenko construction is on the honest side of the
Kapustin--Fidkowski wall~\cite{kapustin-fidkowski}: it computes the Hall
conductance of genuinely chiral states, which no commuting-projector model
realizes---a distinction that becomes Part~V's subject.

\subsection{The SSH chain and its EFT}

\begin{example}[SSH]
\label{ex:ssh}
The Su--Schrieffer--Heeger chain~\cite{ssh-1979} has Bloch Hamiltonian
\[
  H(k;t_1,t_2)=(t_1+t_2\cos k)\,\sigma_x+(t_2\sin k)\,\sigma_y,
  \qquad q(k)=t_1+t_2 e^{ik},
\]
with spectrum $E_\pm(k)=\pm|q(k)|$, gapped iff $|t_1|\ne|t_2|$. It is a class
$\mathit{AIII}$ system in $d=1$, whose table entry (\Cref{tab:tenfold}) is $\Zp$,
the winding number of $q:S^1\to\mathbb C^\times$: $\nu=1$ for $|t_1|<|t_2|$ and
$\nu=0$ for $|t_1|>|t_2|$. The low-energy theory near the gap-closing at
$|t_1|=|t_2|$, $k=\pi$ is a massive $1{+}1$D Dirac fermion, and the winding number
is the sign of the Dirac mass, i.e.\ the deformation class of the Dirac EFT.
Thus $c_{1,\mathit{AIII}}$ sends the SSH phase to its EFT class and the two agree:
Conjecture IV-2 at (L0) is verified in this box. In the condensed formulation the
parameter torus and the momentum circle are replaced by their condensations
$\cond B,\cond{S^1}$ and the winding becomes a value of the family invariant, but
the number is unchanged.
\end{example}

\subsection{Transitions and the spectral boundary map}

A path in the parameter space $B$ that changes the phase label must cross the
gapless discriminant $\Sig_f$. The program organizes this by a relative class:
for a family $f:\cond B\to\Ham_{d,G}$ with gapped locus
$U_f=\cond B\times_{\Ham_{d,G}}\Gap_{d,G}$ and label $\nu_f:U_f\to\pi_0\Phase_{d,G}$,
the obstruction to extending $\nu$ across $\Sig_f=\cond B\setminus U_f$ is the
relative transition charge $\rc\in E^{q+1}(B,U_f)$. We use the boxed slogan
verbatim:
\begin{center}
\fbox{\parbox{0.86\linewidth}{\centering\emph{A topological transition is
crossing the gapless discriminant in the Hamiltonian moduli stack.}}}
\end{center}

\begin{conjecture}[\textbf{IV-3}: relative charge $=$ spectral boundary map]
\label{conj:iv3}
When the invariant $\nu$ is valued in the (condensed) generalized cohomology
represented by $\IPcond_{d,G}$, the relative transition charge $\rc\in
E^{q+1}(B,U_f)$ is the image of $\nu_f$ under the connecting homomorphism of
the cofiber sequence
\[
  \IPcond_{d,G}\text{-cohomology of }U_f
  \ \longrightarrow\
  E^{q+1}(B,U_f)
\]
associated to the pair $(B,U_f)$; that is, the jump of the phase label across
$\Sig_f$ is computed by the boundary map of $\IPcond_{d,G}$. For the SSH chain
(\Cref{ex:ssh}) this recovers the statement that the winding number jumps by
$\pm1$ across $|t_1|=|t_2|$, the linking charge of the gapless point.
\end{conjecture}

Conjecture IV-3 is what makes the spectrum $\IPcond$ do work beyond
classification: it turns the transition calculus into a long-exact-sequence
computation. Its low-dimensional shadow---the SSH winding jump, the Dirac node as
a source of charge---is classical; the conjecture is that this is systematically
the boundary map of $\IPcond$ in every dimension and generalized theory.

\section{Formal verification}
\label{sec:haskell}

The algebraic backbone of the paper (the commutative-monoid structure, its group
completion, and the Bott periodicity of the tenfold-way table) is elementary
enough to be checked by machine, and we do so. The accompanying Haskell package
\texttt{src/\allowbreak lattice-eft-\allowbreak equivalence/} contains four
modules.

\paragraph{\texttt{Monoid.hs}: finitely presented commutative monoids and group
completion.} A finitely presented commutative monoid is represented by generators
and relations; elements are normal forms (multisets of generators modulo the
relations, reduced by a confluent rewriting of the given relations). The
Grothendieck group completion of \Cref{def:groth} is implemented directly:
$\Groth(M)$ is presented as formal differences $[(a,b)]$ with equality decided by
the stable-equality test of \Cref{prop:stable}---search for a witness $e$ with
$a+b'+e=a'+b+e$ over the (finite, for our examples) reachable set. The universal
property (\Cref{thm:groth}(iii)) is exercised by constructing the induced
homomorphism $\bar f$ from a sample $f:M\to A$ and checking
$\bar f\circ\gamma=f$.

\paragraph{\texttt{TenfoldWay.hs}: the periodic table.} The invariant groups
$\mathcal G_{\mathbb C}$ and $\mathcal G_{\mathbb R}$ of \Cref{def:tenfold} are
encoded from the Bott sequences $d_{\mathbb C},d_{\mathbb R}$ via the
antidiagonal formula, and \Cref{tab:tenfold} is regenerated from the code.

\paragraph{\texttt{Main.hs}: demonstrations.} Runs group completion on several
monoids---the free monoid $\Zp_{\ge0}$ (completing to $\Zp$), a product
$\Zp_{\ge0}^2$ (completing to $\Zp^2$), and the non-cancellative toy
$M=\langle a,t\mid t+a=t\rangle$ of \Cref{ex:noncanc}, where it exhibits $[a]=0$
while the multiples $[nt]$ stay distinct, so $\Groth(M)\cong\Zp$---then prints
\Cref{tab:tenfold} and exits with status $0$.

\paragraph{\texttt{Properties.hs}: QuickCheck.} The properties tested correspond
one-to-one to the paper's claims:
\begin{itemize}
\item commutativity, associativity, and unit laws of $\boxtimes$ on monoid
  elements (\Cref{prop:monoid});
\item that $\Groth(M)$ is an abelian group (inverses and associativity) and that
  $\gamma_M$ is injective exactly on the cancellative quotient
  (\Cref{prop:stable});
\item the universal property: for random $f:M\to A$, the induced $\bar f$ is a
  homomorphism and factors $f$ (\Cref{thm:groth});
\item Bott periodicity: $\mathcal G_{\mathbb C}(s,d)=\mathcal G_{\mathbb C}(s,d+2)$,
  $\mathcal G_{\mathbb R}(s,d)=\mathcal G_{\mathbb R}(s,d+8)$, and the antidiagonal
  invariance $\mathcal G(s,d)=\mathcal G(s+1,d+1)$ (\Cref{prop:bott}), together
  with spot checks of the landmark entries of \Cref{tab:tenfold}.
\end{itemize}
The code is verification, not proof: it checks the finite and periodic content of
the elementary claims and guards against off-by-one errors in the table. The
conjectures IV-1 through IV-5 are, by their nature, outside its reach.

\paragraph{Code availability.} The Haskell package is available at
\href{https://github.com/YonedaAI/topological-phases-of-matter}{github.com/YonedaAI/topological-phases-of-matter},
under \texttt{src/\allowbreak lattice-eft-\allowbreak equivalence/}. It builds
with GHC and its QuickCheck suite runs the properties listed above.

\section{Discussion}
\label{sec:discussion}

\paragraph{What is honestly established.} The symmetric-monoidal structure on
$\pi_0$ and its group completion (\Cref{prop:monoid,thm:groth,prop:stable}) are
theorems; the reading of ancilla stabilization as the group-completion
stabilization element (\Cref{rem:stmeaning}) is, we think, the clarifying point
of the paper, and it is elementary. The recognition of the invertible sector as a
connective spectrum (\Cref{thm:recognition}) is a theorem \emph{conditional} on
$\Phase_{d,G}$ being a symmetric-monoidal $\infty$-groupoid, and its output has a
rigorous carrier in Kubota's $\Omega$-spectrum (\Cref{thm:kubota}). The
free-fermion table and its periodicities (\Cref{prop:bott}) are classical and
machine-checked. The $d=1$ comparison (\Cref{thm:d1}) is a genuine theorem at
level (L0).

\paragraph{What is conjectural, and why we did not hide it.} The five conjectures
IV-1 through IV-5 are the substance of the program at this floor, and none is a
theorem. The temptation in a paper like this is to state IV-2 as a ``theorem
under assumptions'' and bury the assumptions; we have instead separated the three
levels of equivalence (\Cref{subsec:levels}) precisely so that the true logical
status is visible: (L0) is a theorem in two regimes, (L1) has a target but no
identification, (L2) is not even statable without the condensed refinement IV-1.

\paragraph{The EFT hypotheses are real.} Every positive statement is under
short-range entanglement. A lattice system is not automatically an invertible
field theory; the assertion that it flows to one is the SRE hypothesis, and it
fails for intrinsic topological order. The Freed--Hopkins classification is proved
\emph{under} reflection positivity and the EFT axioms~\cite{freed-hopkins}; we
inherit those hypotheses and do not claim more.

\paragraph{Stabilization collapses distinctions, on purpose.}
\Cref{prop:stable} and \Cref{ex:noncanc} show that group completion (and the
$(-)^{\st}$ built into $\Phase_{d,G}$) can send a nonzero phase to zero once
enough ancillas are added. This is intended: it is what makes ``phase'' a stable
notion, but it means $\IPcond$ sees only stable, invertible data. Unstable or
fragile distinctions are invisible to it by construction.

\paragraph{Noninvertible order is out of scope.} The whole paper lives in the
invertible sector. Intrinsic topological order (anyons, modular tensor categories
in $2{+}1$D) has no inverse under $\boxtimes$ and no place in a spectrum; a
condensed treatment of it would need a condensed \emph{higher} stack of phases and
defects, not a $K$-theory spectrum, and is deferred entirely
(\cite{paperSynthesis}, global pitfalls). The De\,Nittis(--Rendel) state-space
programme~\cite{denittis-states,denittis-rendel-weyl} and the moduli-space work of
Hsin--Wang~\cite{hsin-wang-moduli} are the nearest ordinary-topology precedents;
our contribution over them is organizational: placing the same invariants inside
a condensed environment where families, disorder, and completions cohabit, not a
new invariant, as the program concedes throughout.

\paragraph{Undecidability bounds the renormalization story.} Conjecture IV-5 is
asserted only on the SRE stratum and only as an existence statement, because the
undecidability of the spectral gap~\cite{cpw-undecidable} forbids any general
procedure certifying that a coarse-graining flows to a gapped fixed point. The
same wall constrains Part~III; the program does not pretend to see past it.

\section{Conclusion}
\label{sec:conclusion}

We have built the top floor of the condensed tower for topological phases and
stated precisely what it would take for the microscopic and field-theoretic
classifications to coincide. The provable content is the algebra of stacking:
$\pi_0$ is a commutative monoid, its group completion is governed by a universal
property, and the stabilization element of that completion is exactly ancilla
stabilization. The invertible sector is a connective spectrum $\IPcond_{d,G}$ with
a rigorous carrier in Kubota's $\Omega$-spectrum. Against the Freed--Hopkins
bordism target we posed the comparison conjecture IV-2 at three levels of
strength and marshalled the evidence that pins it down where it can be pinned:
Ogata's completeness makes the $\pi_0$ comparison a theorem in $d=1$, and the
tenfold-way $K$-theory, with its Bott periodicities, which we verified by
machine, realizes it for free fermions in every dimension. The remaining
conjectures locate the missing analysis: solidification as the EFT completion
(IV-4, IV-5), the condensed refinement of the spectrum (IV-1), and the transition
charge as a spectral boundary map (IV-3).

The honest summary is that the equivalence between lattice models and effective
field theories is, at $\pi_0$ and under EFT hypotheses, either a theorem
(Ogata, free fermions) or a well-posed and evidence-constrained conjecture, and
that its spectrum- and family-level forms are exactly what the condensed
refinement is for. Part~V takes the baton with the realizability question (which
of these classes are built by actual lattices), and Part~VI assembles the tower
into the program's Master Conjecture, of which our IV-2 is the top.

\bibliographystyle{plain}
\bibliography{references}

\end{document}
